\section[Skeleton]{The skeleton}

\blockframe{1}{The skeleton}{From an empty file to a structured PDF}

\begin{frame}
    \frametitle{Schedule}
    \vfill
    \begin{roadmap}{1}
        \roadmapresults
    \end{roadmap}
    \vfill
    \note{%
        Second showing.
        Point at the thick part of the line: ``this one, in the next 20 minutes.''
    }
\end{frame}

% ---------------------------------------------------------------------
%  Anatomy of a command
% ---------------------------------------------------------------------
% The \tikzmark's sit between the parts of the command (c0 to c4).
% \annotatebelow/\annotateabove then draw a bracket between two marks.
% Overlays 1-4 show the command without, overlay 5 on with the optional
% argument; the marks have the same names in both, so the brackets follow.
\begin{frame}[fragile]
    \frametitle{Commands}
    \begin{center}
        \LARGE
        \begin{onlyenv}<1-4| handout:1>
            \mintinline[escapeinside=||]{latex}{|\tikzmark{c0}|\|\tikzmark{c1}|documentclass|\tikzmark{c2}|{article}|\tikzmark{c4}|}
        \end{onlyenv}
        \begin{onlyenv}<5-| handout:2>
            \mintinline[escapeinside=||]{latex}{|\tikzmark{c0}|\|\tikzmark{c1}|documentclass|\tikzmark{c2}|[a4paper]|\tikzmark{c3}|{article}|\tikzmark{c4}|}
        \end{onlyenv}
    \end{center}
    \begin{tikzpicture}[remember picture, overlay]
        \only<2-| handout:1,2>{\annotatebelow[-1.2cm]{c0}{c1}{3mm}{Command\\delimiter}}
        \only<3-| handout:1,2>{\annotatebelow{c1}{c2}{3mm}{Command\\name}}
        \only<4| handout:1>{\annotatebelow{c2}{c4}{3mm}{Mandatory\\argument}}
        \only<5-| handout:2>{\annotatebelow{c3}{c4}{3mm}{Mandatory\\argument}}
        \only<5-| handout:2>{\annotateabove{c2}{c3}{3mm}{Optional argument}}
    \end{tikzpicture}
    \note{%
        Every command today has this shape: backslash, letters, arguments.
        Go slowly here and do not apologise for it.
        Three minutes here saves twenty minutes of confusion later:
        from now on they can \emph{parse} what they see, even when they do not know what it does.
        \begin{itemize}
            \item Say out loud: a non-letter ends a command name.
                That is why \texttt{\textbackslash LaTeX2e} does not work,
                and why they will one day need \texttt{\textbackslash LaTeX\{\} is nice}.
                Plant it, do not dwell.
        \end{itemize}
    }
\end{frame}

% ---------------------------------------------------------------------
%  Anatomy of an environment
% ---------------------------------------------------------------------
% Marks: n1l/n1r and n2l/n2r around the two environment names,
% b0 and e0 at the start of the \begin and \end lines.
\begin{frame}[fragile]
    \frametitle{Environments}
    \begin{center}
        \begin{minipage}{4.6cm}
            \begin{minted}[fontsize=\LARGE, escapeinside=||]{latex}
                |\tikzmark{b0}|\begin{|\tikzmark{n1l}|center|\tikzmark{n1r}|}
                    |\textrm{\textit{Content}}|
                |\tikzmark{e0}|\end{|\tikzmark{n2l}|center|\tikzmark{n2r}|}
            \end{minted}
        \end{minipage}
    \end{center}
    \begin{tikzpicture}[remember picture, overlay]
        % Overlay 2: a dashed box around each name, both joined to one label.
        \only<2->{%
            \node[annotation, draw, dashed, inner sep=1pt,
                fit={([yshift=-1.2mm]pic cs:n1l) ([yshift=4.2mm]pic cs:n1r)}] (name1) {};
            \node[annotation, draw, dashed, inner sep=1pt,
                fit={([yshift=-1.2mm]pic cs:n2l) ([yshift=4.2mm]pic cs:n2r)}] (name2) {};
            \node[annotation, anchor=west, font=\normalsize] (namelabel)
                at ($(name1.east)!.5!(name2.east) + (3cm,0)$) {Same name twice};
            \draw[annotation] (name1.north) -- ++(0,2mm) -| (namelabel.north);
            \draw[annotation] (name2.south) -- ++(0,-2mm) -| (namelabel.south);
        }
        % Overlay 3: a bracket on the left, from below \begin to above \end.
        \only<3->{%
            \coordinate (top) at ([xshift=-4mm, yshift=-2.2mm]pic cs:b0);
            \coordinate (bottom) at ([xshift=-4mm, yshift=5mm]pic cs:e0);
            \draw[annotation] ([xshift=\annoarm]top) -- (top) -- (bottom) -- ([xshift=\annoarm]bottom);
            \node[annotation, anchor=east, align=right, font=\normalsize]
                at ($(top)!.5!(bottom) + (-2mm,0)$) {Everything in here\\is inside};
        }
    \end{tikzpicture}
    \note{%
        \begin{itemize}
            \item A command acts \emph{here}; an environment acts on \emph{all of this}.
            \item First the shape, then the names, then the contents.
                It makes the ``wrapping'' idea visual.
            \item Promise them now: a mismatched \texttt{\textbackslash begin}/\texttt{\textbackslash end}
                is the error they will hit most often today.
                When it happens in block 2 and you can say ``told you'',
                they trust the next prediction too.
        \end{itemize}
    }
\end{frame}

\begin{frame}[fragile]
    \frametitle{Reserved Characters}
    \begin{center}
        \renewcommand{\arraystretch}{1.3}
        \begin{tabular}{@{}>{\ttfamily}r l l@{}}
            \toprule
            \textnormal{\textbf{Symbol}} & \textbf{Description} & \textbf{Substitute} \\
            \midrule
            \%                  & comment                    & \code|\%| \\
            \{ \}               & grouping                   & \code|\{ \}| \\
            \$                  & math on and off           & \code|\$| \\
            \&                  & table column separator     & \code|\&| \\
            \#                  & argument placeholder       & \code|\#| \\
            \_ \^{}             & subscript, superscript     & \code|\_| \code|\^{}| \\
            \textasciitilde     & non-breaking space         & \code|\textasciitilde| \\
            \textbackslash      & start of a command         & \code|\textbackslash| \\
            \bottomrule
        \end{tabular}
    \end{center}
    \note{%
        \begin{itemize}
            \item Do not read the table out.
                Point at \texttt{\%}, say the warning (it kills the rest of the line), move on.
            \item The third column answers the question people actually ask:
                ``how do I type a percent sign?''
            \item The backslash: \texttt{\textbackslash\textbackslash} is \emph{not} how you escape it.
                That is a line break; \texttt{\textbackslash textbackslash} is the character.
            \item The tilde comes back in block 2 and 4, as the space in \texttt{Fig.\textasciitilde\textbackslash ref}.
        \end{itemize}
    }
\end{frame}

% Three states of the same document. The lines that are new in each state
% get a coloured background (highlightlines).
\begin{frame}
    \frametitle{Minimal Document}
    \begin{center}
        \only<1| handout:1>{\tcbinputlst[][width=.8\textwidth]{listings/01-minimal-1.tex}}%
        \only<2| handout:2>{\tcbinputlst[highlightlines={3-5,8}][width=.8\textwidth]{listings/01-minimal-2.tex}}%
        \only<3| handout:3>{\tcbinputlst[fontsize=\footnotesize, highlightlines={10-14}][width=.8\textwidth]{listings/01-minimal-3.tex}}%
    \end{center}
    \note{%
        Three clicks: it compiles, it has a title, it has a shape.
        Nothing else.
        \begin{itemize}
            \item \texttt{\textbackslash documentclass} is \emph{outside} \texttt{\textbackslash begin\{document\}}.
                Someone always asks.
                That region is the preamble, and it comes back in block 2.
        \end{itemize}
    }
\end{frame}

\begin{frame}
    \frametitle{Sections}
    \begin{columns}[T, onlytextwidth]
        \begin{column}{.56\textwidth}
            \only<1| handout:1>{\tcbinputlst[fontsize=\footnotesize]{listings/01-sections.tex}}%
            \only<2-| handout:2>{\tcbinputlst[fontsize=\footnotesize, highlightlines={5,8,10}]{listings/01-sections.tex}}%
        \end{column}
        \begin{column}{.38\textwidth}
            % trim = left bottom right top, in bp: the heading area of page 1
            \includegraphics[width=\dimexpr\textwidth-0.8pt\relax, frame, trim=110 606 265 118, clip]{renders/01-sections.pdf}

            \bigskip
            \uncover<3-| handout:2>{\Large Numbering is automatic}
        \end{column}
    \end{columns}
    \note{%
        \begin{itemize}
            \item Hold on that last line.
                The numbering is the first thing all afternoon that \LaTeX{} did \emph{for} them,
                not \emph{to} them.
                It sets up the table of contents on the next slide.
            \item Ask the room: ``what happens to those numbers if I move Theory above Introduction?''
                Let someone say it.
        \end{itemize}
    }
\end{frame}

\begin{frame}
    \frametitle{Table of Contents}
    \begin{columns}[T, onlytextwidth]
        \begin{column}{.56\textwidth}
            \tcbinputlst[fontsize=\footnotesize, highlightlines={4}]{listings/01-toc.tex}
        \end{column}
        \begin{column}{.4\textwidth}
            % trim = left bottom right top, in bp: the Contents block of page 1
            \includegraphics[width=\dimexpr\textwidth-0.8pt\relax, frame, trim=110 607 110 118, clip]{renders/01-toc.pdf}

            \medskip
            \uncover<2->{%
                Built from your \texttt{\textbackslash section} commands
            }
        \end{column}
    \end{columns}
    \note{%
        \textbf{The key slide of block 1.}
        Nobody typed that list, and it stays right when they rename or reorder sections.
        Describe the structure once, and \LaTeX{} keeps the rest up to date.
        If they remember one thing from the first half hour, make it this one.
        \begin{itemize}
            \item Ask what else could be generated from structure they already wrote.
                Someone will guess page numbers, cross-references, the bibliography.
                ``Yes, all of that, before half past five.''
            \item Deliberately \emph{not} teaching two compilation passes here.
                Overleaf and \texttt{latexmk} handle it.
                The one case that still bites (a misspelled citation key) comes in block 4.
        \end{itemize}
    }
\end{frame}

% ---------------------------------------------------------------------
%  Text styling: the commands, then an example, then \emph
% ---------------------------------------------------------------------
% One row per click. \onslide is used as a switch inside the cells
% because \code (verbatim) cannot sit inside the argument of \uncover.
\begin{frame}[fragile]
    \frametitle{Text Styling}
    \begin{center}
        \Large
        \renewcommand{\arraystretch}{1.5}
        \begin{tabular}{@{}l@{\hspace{3em}}l@{}}
            \toprule
            \textbf{Command} & \textbf{Output} \\
            \midrule
            \onslide<1->\code|\textbf{bold}|          & \onslide<1->\textbf{bold} \\
            \onslide<2->\code|\textit{italic}|        & \onslide<2->\textit{italic} \\
            \onslide<3->\code|\underline{underlined}| & \onslide<3->\underline{underlined} \\
            \onslide<4->\code|\texttt{typewriter}|    & \onslide<4->\texttt{typewriter} \\
            \onslide<5->\code|\textsf{sans-serif}|    & \onslide<5->\textsf{sans-serif} \\
            \bottomrule
        \end{tabular}
    \end{center}
    \note{%
        Click through quickly; they all have the same shape as \texttt{\textbackslash documentclass}.
        \begin{itemize}
            \item \texttt{\textbackslash underline} does not break across lines.
                For longer underlined text there is the \texttt{ulem} package (handout).
            \item \texttt{\textbackslash texttt} is the one they will use most after bold and italic:
                file names, commands, code.
            \item More commands (\texttt{\textbackslash textsc}, \texttt{\textbackslash textsl}, \ldots) are in the handout.
        \end{itemize}
    }
\end{frame}

\begin{frame}
    \frametitle{Text Styling: Example}
    \vfill
    \tcbinputlst[fontsize=\footnotesize]{listings/01-styling.tex}
    \uncover<2->{\outputrender{renders/01-styling.pdf}}
    \vfill
    \note{%
        The same commands in a real sentence.
        TODO 1.5 asks for exactly these three.
    }
\end{frame}

\begin{frame}[fragile]
    \frametitle{Emphasis}
    \vfill
    \tcbinputlst[fontsize=\footnotesize]{listings/01-emph.tex}
    \uncover<2->{\outputrender{renders/01-emph.pdf}}
    \vfill
    \begin{uncoverenv}<3->
        \begin{center}
            \color{black!60}
            How \code|\emph| looks is set in the preamble $\rightarrow$ block 2
        \end{center}
    \end{uncoverenv}
    \note{%
        \texttt{\textbackslash emph} marks text as emphasised compared to what is around it.
        \begin{itemize}
            \item In upright text that means italic, in italic text upright.
            \item Do not go further here.
                Block 2 changes how emphasis looks for the whole document, with one line in the preamble.
            \item If someone asks why it does not flip on their beamer slides:
                beamer redefines \texttt{\textbackslash emph} as ``always italic''.
        \end{itemize}
    }
\end{frame}

\begin{frame}
    \frametitle{Summary}
    \tcbinputlst[fontsize=\footnotesize]{listings/01-summary.tex}
    \vfill
    \begin{center}
        \color{black!60}
        Also in \texttt{checkpoints/checkpoint-01.tex}
    \end{center}
    \note{%
        A reference slide, not a teaching one.
        Leave it up while they work if you like.
    }
\end{frame}

\begin{frame}
    \frametitle{Exercise}
    \handson{1}{1.1--1.5}{5--10}{checkpoint-01.tex}{%
        Open \texttt{workshop.tex} and work down the TODOs.
        \begin{itemize}
            \normalsize\color{black!60}
            \setbeamercolor{itemize item}{fg=black!60}
            \item Compiling document
            \item Title block
            \item Sections
            \item Table of contents
            \item Styled words
        \end{itemize}
    }
    \note{%
        \textbf{5 to 10 minutes.}
        Set a timer and show it.
        \begin{itemize}
            \item Walk the room.
                Only answer from the front if three people have the same problem;
                then stop everyone and fix it once, out loud.
            \item The two failures you will see:
                \texttt{\textbackslash begin\{document\}} without \texttt{\textbackslash end\{document\}},
                and a smart quote or long dash pasted in from somewhere.
            \item Two minutes before you stop: ``if you are behind, grab checkpoint-01.''
                Say it loudly and without judgement.
            \item \textbf{Break after this slide}, 5 minutes.
                Say the actual restart time, and start then.
        \end{itemize}
    }
\end{frame}
